% Rozdział 5 — Badania eksperymentalne (budżet: 13 str.)
% 5.1 Źródło: emisja-magnes-pomiary/emisja_tresc.tex (zaczyna się od \section).
% Fragment używa \figdir i \datadir — ustawiane tutaj względem wspolne/.
\chapter{Badania eksperymentalne}
\label{ch:badania}

W~rozdziale przedstawiono wyniki pomiarów emisji polowej prototypu
z~magnesem ogniskującym i~bez niego oraz widma promieniowania X
wzbudzanego wiązką elektronów w~próbkach o~różnym składzie. Stanowisko
i~metody opisano w~\rozdzref{ch:stanowisko}.

\renewcommand{\figdir}{../emisja-magnes-pomiary/figures}
\renewcommand{\datadir}{../emisja-magnes-pomiary/dane}
% =====================================================================
%  Emisja polowa: sesja z magnesem vs. sesje bez magnesu
%  17.11.2026: U -> \UB, intensywność I -> jasność L, e -> \Delta, r -> r_P,
%  sesje nazywane „bez magnesu”/„z magnesem” (nie B/M: B = bramka),
%  bez nazw plików roboczych; metoda analizy zdjęć przeniesiona do rozdz. 4.
%  Fragment do wstawienia w rozdział 5 (\input{...}).
%  Wymaga: graphicx, subcaption, pgfplots (compat>=1.17), booktabs, siunitx
%  Ścieżki: figures/ (zdjęcia), dane/ (CSV z pomiarów i analizy obrazu)
% =====================================================================
\providecommand{\sesjaM}{z~magnesem}
\providecommand{\sesjaB}{bez magnesu}
\providecommand{\figdir}{figures}
\providecommand{\datadir}{dane}
\pgfplotsset{/pgf/number format/use comma, /pgf/number format/1000 sep={}}
\providecommand{\TODO}[1]{\textcolor{red}{[TODO: #1]}}

\section{Wpływ magnesu ogniskującego na emisję polową}
\label{sec:emisja-magnes}

\subsection{Przebieg pomiarów i~charakterystyki prądowo-napięciowe}

Porównano dwie sesje pomiarowe, w~których przy stałym napięciu $\UT$
zwiększano napięcie $\UB$ co \SI{100}{V} i~rejestrowano prąd emisji
$\Ie$ oraz zdjęcie ekranu luminescencyjnego (\rozdzref{ch:stanowisko}):
\begin{itemize}
  \item \textbf{sesja bez magnesu} -- $\UT=\SI{2000}{V}$,
        $\UB=\SIrange{500}{1500}{V}$, prąd w~zakresie
        \SIrange{2}{22}{\micro A}, zdjęcia dla
        $\UB=\SIrange{700}{1100}{V}$; przed tą sesją źródło odgazowano
        przez wygrzanie w~temperaturze ok. \SI{200}{\celsius},
  \item \textbf{sesja z~magnesem} -- $\UT=\SI{1500}{V}$,
        $\UB=\SIrange{500}{1700}{V}$, prąd w~zakresie
        \SIrange{14}{200}{\micro A}, zdjęcia dla
        $\UB=\SIrange{500}{1600}{V}$; źródło nie było przed tą sesją
        ponownie odgazowane.
\end{itemize}
Dodatkowo wykorzystano serię~2 bez magnesu (10~zdjęć z~odczytami prądu
\SIrange{1}{400}{\micro A}), która posłużyła do sprawdzenia związku jasności
ekranu z~prądem. Wyniki zebrano w~tab.~\ref{tab:emisja-iv}.

\begin{table}[htbp]
  \centering
  \caption{Prąd emisji w~funkcji napięcia w~obu sesjach}
  \label{tab:emisja-iv}
  \begin{tabular}{@{}r r r@{}}
    \toprule
    $\UB$ [V] & $\Ie$ \sesjaB\ [µA] & $\Ie$ \sesjaM\ [µA] \\
    \midrule
     500 & 0    & 14 \\
     600 & 0    & 16 \\
     700 & 0    & 14 \\
     800 & 2    & 28 \\
     900 & 3,1  & 45 \\
    1000 & 4    & 69 \\
    1100 & 6    & 118 \\
    1200 & 8    & 180 \\
    1300 & 9    & 190 \\
    1400 & 11,2 & 190 \\
    1500 & 22   & 193 \\
    1600 & --   & 200 \\
    1700 & --   & 200 \\
    \midrule
    $\UT$ & \SI{2000}{V} & \SI{1500}{V} \\
    \bottomrule
  \end{tabular}
\end{table}

Uzyskane charakterystyki $\Ie(\UB)$ przedstawiono na rys.~\ref{fig:emisja-iv}.
W~sesji bez magnesu prąd pojawia się skokowo przy \SI{800}{V} i~rośnie niemal
liniowo do \SI{1400}{V}; punkt \SI{1500}{V} (\SI{22}{\micro A}, podwojenie
prądu) odbiega od trendu i~może wskazywać na początek wyładowania
lub niestabilności emisji. W~sesji z~magnesem dla $\UB\leq\SI{700}{V}$ płynie
stały prąd ok. \SI{15}{\micro A}, który należy traktować jako prąd
upływu (tło), następnie prąd rośnie wykładniczo do \SI{180}{\micro A}
przy \SI{1200}{V}, po czym nasyca się na poziomie
\SIrange{190}{200}{\micro A}. Przy tym samym napięciu $\UB$ prąd w~sesji
z~magnesem jest ok. 10--20 razy większy niż bez magnesu (np. \SI{69}{\micro A} wobec
\SI{4}{\micro A} przy \SI{1000}{V}), mimo niższego $\UT$.

Interpretację tych wyników utrudnia to, że sesje nie różnią się wyłącznie obecnością magnesu. Przed sesją bez
magnesu źródło odgazowano, a~przed sesją z~magnesem nie powtórzono tej
procedury. Gazy zaadsorbowane na powierzchni emitera i~elementów
źródła zmieniają pracę wyjścia emitera i~mogą być desorbowane pod
wpływem wiązki, lokalnie zwiększając ciśnienie (\pref{sec:emisja}).
Prawdopodobnie dlatego w~sesji z~magnesem obserwuje się:
\begin{itemize}
  \item stały prąd ok. \SI{15}{\micro A} już przy \SI{500}{V}, przed
        właściwym progiem emisji (upływ lub prąd jonowy),
  \item nasycenie prądu powyżej \SI{1200}{V},
  \item rozległą, rozmytą poświatę wokół plamki, która może być
        częściowo świeceniem gazu resztkowego.
\end{itemize}
Wyższy prąd w~sesji z~magnesem może wynikać z~dwóch przyczyn, których na
podstawie tych pomiarów nie da się rozdzielić: ogniskowania wiązki
przez magnes (więcej elektronów trafia w~target zamiast w~elektrody
i~ścianki źródła) oraz innego stanu powierzchni emitera i~ciśnienia
w~nieodgazowanym źródle, łącznie z~udziałem prądu upływu. Dlatego
wpływ magnesu oceniono przede wszystkim na podstawie kształtu
i~skupienia plamki.

\begin{figure}[htbp]
  \centering
  \begin{tikzpicture}
    \begin{semilogyaxis}[
      width=0.85\textwidth, height=7cm,
      xlabel={napięcie $\UB$ [V]}, ylabel={prąd emisji $\Ie$ [µA]},
      xmin=450, xmax=1750, ymin=1, ymax=500, xtick distance=200,
      grid=major, grid style={gray!25},
      legend pos=south east, legend cell align=left,
      every axis plot/.append style={thick, mark size=2pt},
      unbounded coords=jump,
    ]
      \addplot[red!60!black, mark=triangle*] table[x=U,y=I_uA,col sep=comma]{\datadir/iv_bez.csv};
      \addlegendentry{\sesjaB\ ($\UT=\SI{2000}{V}$)}
      \addplot[blue!70!black, mark=*] table[x=U,y=I_uA,col sep=comma]{\datadir/iv_mag.csv};
      \addlegendentry{\sesjaM\ ($\UT=\SI{1500}{V}$)}
      \draw[gray, dashed] (axis cs:450,15) -- (axis cs:1750,15)
        node[pos=0.2, below, font=\scriptsize] {upływ ok. \SI{15}{\micro A}};
    \end{semilogyaxis}
  \end{tikzpicture}
  \caption{Charakterystyki prądowo-napięciowe emisji w~obu sesjach
    (skala logarytmiczna; punkty z~$\Ie=0$ pominięto)}
  \label{fig:emisja-iv}
\end{figure}


\subsection{Obraz wiązki na ekranie luminescencyjnym}

Jasność świecenia ekranu $L$ wyznaczono ze zdjęć metodą opisaną
w~\pref{sec:analiza-zdjec}. Na rys.~\ref{fig:emisja-I} zestawiono ją
z~prądem emisji w~funkcji napięcia. W~sesji bez magnesu jasność ekranu
rośnie razem z~prądem ($r_\mathrm{P}=0{,}99$ dla \SIrange{800}{1100}{V}).
W~sesji z~magnesem $r_\mathrm{P}=0{,}95$ dla \SIrange{800}{1200}{V}, a~powyżej
\SI{1200}{V} zarówno prąd, jak i~jasność przestają rosnąć -- nasycenie
prądu jest więc rzeczywiste, a~nie wynika tylko z~prześwietlenia
kamery. W~sesji z~magnesem słaba plamka jest widoczna już
przy \SIrange{600}{700}{V}, gdy mierzony prąd nie odbiega jeszcze od
poziomu upływu -- magnes skupia niewielki prąd emisji na małym
obszarze ekranu, przez co staje się on widoczny wcześniej niż bez
magnesu (próg \SI{800}{V}). Dla serii~2 \sesjaB\ korelacja jasności
z~odczytem prądu wynosi $r_\mathrm{P}=0{,}98$ (rys.~\ref{fig:emisja-IvsIe}).

\begin{figure}[htbp]
  \centering
  \begin{tikzpicture}
    \begin{axis}[
      width=0.8\textwidth, height=6.5cm,
      xlabel={napięcie $\UB$ [V]}, ylabel={jasność $L$ [j.u.]},
      xmin=450, xmax=1650, ymin=0, ymax=3.5, xtick distance=200,
      grid=major, grid style={gray!25},
      legend pos=north west, legend cell align=left,
      every axis plot/.append style={thick, mark size=2pt},
      axis y line*=left,
    ]
      \addplot[blue!70!black, mark=*] table[x=U,y=I,col sep=comma]{\datadir/mg.csv};
      \addlegendentry{jasność, \sesjaM}
      \addplot[red!60!black, mark=triangle*] table[x=U,y=I,col sep=comma]{\datadir/s1.csv};
      \addlegendentry{jasność, \sesjaB}
      \addlegendimage{blue!70!black, dashed, mark=o}
      \addlegendentry{prąd, \sesjaM\ (oś prawa)}
    \end{axis}
    \begin{axis}[
      width=0.8\textwidth, height=6.5cm,
      xmin=450, xmax=1650, ymin=0, ymax=220,
      axis y line*=right, axis x line=none,
      ylabel={$\Ie$ \sesjaM\ [µA]},
    ]
      \addplot[blue!70!black, dashed, mark=o, thick] table[x=U,y=I_uA,col sep=comma]{\datadir/iv_mag.csv};
    \end{axis}
  \end{tikzpicture}
  \caption{Jasność świecenia ekranu (oś lewa) i~prąd emisji
    w~sesji z~magnesem (oś prawa) w~funkcji napięcia}
  \label{fig:emisja-I}
\end{figure}

\begin{figure}[htbp]
  \centering
  \begin{tikzpicture}
    \begin{loglogaxis}[
      width=0.65\textwidth, height=5.5cm,
      xtick={30,50,100,200,400}, xticklabels={30,50,100,200,400},
      log ticks with fixed point,
      xlabel={prąd emisji $\Ie$ [µA] (seria 2, \sesjaB)},
      ylabel={jasność $L$ [j.u.]},
      grid=major, grid style={gray!25},
      legend pos=north west, legend cell align=left,
    ]
      \addplot[only marks, mark=square*, orange!85!black]
        table[x=etykieta, y=I, col sep=comma,
              restrict expr to domain={\thisrow{I}}{0.01:100}]{\datadir/s2.csv};
      \addlegendentry{pomiar}
      \addplot[domain=30:420, samples=20, black, dashed] {exp(-6.886)*x^1.371};
      \addlegendentry{$L\propto \Ie^{1{,}37}$}
    \end{loglogaxis}
  \end{tikzpicture}
  \caption{Jasność świecenia a~odczyt prądu w~serii~2 \sesjaB}
  \label{fig:emisja-IvsIe}
\end{figure}

Wpływ magnesu na kształt wiązki oceniono na podstawie zdjęć obu sesji
przy tych samych napięciach \SIrange{800}{1100}{V}, które zestawiono na rys.~\ref{fig:emisja-sbs}, a~parametry plamki w~tab.~\ref{tab:emisja}.
\begin{itemize}
  \item \textbf{Bez magnesu} świeci duży fragment okna ekranu o~ostrych
        krawędziach -- kształt plamki odtwarza geometrię otworu, czyli
        wiązka jest szeroka i~równomiernie oświetla okno.
  \item \textbf{Z~magnesem} świeci zwarty, wydłużony obszar w~części okna,
        otoczony rozległą niebieską poświatą z~promienistymi smugami.
\end{itemize}
Obraz z~sesji z~magnesem jest zgodny z~działaniem magnesu jako soczewki
ogniskującej: wiązka zostaje ściśnięta do mniejszego obszaru, zamiast
oświetlać całe okno, dzięki czemu plamka jest widoczna już przy
mniejszym prądzie. Poświata i~smugi mogą pochodzić częściowo
z~rozpraszania światła w~optyce kamery (prześwietlony rdzeń),
a~częściowo z~elektronów poza głównym ogniskiem lub ze świecenia
gazu resztkowego w~nieodgazowanym źródle. Ze względu na brak
skali na zdjęciach i~inne kadrowanie nie można podać średnicy plamki
w~mikrometrach.

\begin{table}[htbp]
  \centering
  \caption{Parametry emisji i~plamki w~obu sesjach}
  \label{tab:emisja}
  \begin{tabular}{@{}lccc@{}}
    \toprule
    Parametr & \sesjaM & \sesjaB & seria 2 (\sesjaB) \\
    \midrule
    Odgazowanie źródła & nie & tak (\SI{200}{\celsius}) & -- \\
    Napięcie $\UT$ & \SI{1500}{V} & \SI{2000}{V} & -- \\
    Próg emisji (prąd) & \SI{800}{V} & \SI{800}{V} & -- \\
    Pierwsza widoczna plamka & \SI{600}{V} & \SI{800}{V} & -- \\
    Prąd maksymalny & \SI{200}{\micro A} & \SI{22}{\micro A} & \SI{400}{\micro A} \\
    Maks. jasność $L$ [j.u.] & 3,19 & 1,67 & 2,89 \\
    Maks. pole plamki [\% kadru] & 3,98 & 2,07 & 3,91 \\
    \bottomrule
  \end{tabular}
\end{table}

% ---------------- zestawienie side-by-side ----------------
\newcommand{\sbs}[2]{%
  \begin{subfigure}[t]{0.28\textwidth}\centering
    \includegraphics[width=\linewidth]{\figdir/#1}\caption{#2}
  \end{subfigure}}

\begin{figure}[p]
  \centering
  \makebox[0.28\textwidth]{\textbf{\sesjaM}}\hspace{1.5em}\makebox[0.28\textwidth]{\textbf{\sesjaB}}\par\medskip
  \sbs{mg_03.jpg}{\SI{800}{V}, \SI{28}{\micro A}}\hspace{1.5em}
  \sbs{s1_02.jpg}{\SI{800}{V}, \SI{2}{\micro A}}\par\medskip
  \sbs{mg_04.jpg}{\SI{900}{V}, \SI{45}{\micro A}}\hspace{1.5em}
  \sbs{s1_03.jpg}{\SI{900}{V}, \SI{3,1}{\micro A}}\par\medskip
  \sbs{mg_05.jpg}{\SI{1000}{V}, \SI{69}{\micro A}}\hspace{1.5em}
  \sbs{s1_04.jpg}{\SI{1000}{V}, \SI{4}{\micro A}}\par\medskip
  \sbs{mg_06.jpg}{\SI{1100}{V}, \SI{118}{\micro A}}\hspace{1.5em}
  \sbs{s1_05.jpg}{\SI{1100}{V}, \SI{6}{\micro A}}
  \caption{Plamka luminescencji \sesjaM\ (lewa kolumna) i~\sesjaB\ (prawa
    kolumna) przy tych samych napięciach -- zbliżenia okna ekranu.
    Z~magnesem świeci zwarty obszar z~poświatą, bez magnesu -- szeroki
    fragment okna o~ostrych krawędziach}
  \label{fig:emisja-sbs}
\end{figure}

\begin{figure}[htbp]
  \centering
  \begin{subfigure}[t]{0.28\textwidth}\centering
    \includegraphics[width=\linewidth]{\figdir/mg_02.jpg}\caption{\SI{700}{V}}\end{subfigure}\hspace{1.5em}
  \begin{subfigure}[t]{0.28\textwidth}\centering
    \includegraphics[width=\linewidth]{\figdir/mg_09.jpg}\caption{\SI{1400}{V}}\end{subfigure}\hspace{1.5em}
  \begin{subfigure}[t]{0.28\textwidth}\centering
    \includegraphics[width=\linewidth]{\figdir/mg_11.jpg}\caption{\SI{1600}{V}}\end{subfigure}
  \caption{Sesja \sesjaM: plamka widoczna przy prądzie na poziomie upływu
    (\SI{700}{V}) oraz w~zakresie nasycenia prądu (\SIrange{1400}{1600}{V})}
  \label{fig:emisja-mg-skrajne}
\end{figure}

\begin{figure}[htbp]
  \centering
  \begin{subfigure}[t]{0.32\textwidth}\centering
    \includegraphics[width=\linewidth]{\figdir/mg_pelny.jpg}\caption{\sesjaM, \SI{1400}{V}}\end{subfigure}\hfill
  \begin{subfigure}[t]{0.32\textwidth}\centering
    \includegraphics[width=\linewidth]{\figdir/s1_pelny.jpg}\caption{\sesjaB, \SI{1100}{V}}\end{subfigure}\hfill
  \begin{subfigure}[t]{0.32\textwidth}\centering
    \includegraphics[width=\linewidth]{\figdir/s2_pelny.jpg}\caption{\sesjaB, seria 2}\end{subfigure}
  \caption{Pełne kadry przy maksymalnej emisji w~każdej serii}
  \label{fig:emisja-pelne}
\end{figure}

Z~pomiarów emisji polowej wynikają następujące wnioski:
\begin{enumerate}
  \item W~obu sesjach próg emisji wynosi ok. \SI{800}{V}.
  \item Z~magnesem plamka jest zwarta i~widoczna już przy
        \SI{600}{V}--\SI{700}{V}, mimo że prąd nie odbiega jeszcze od
        poziomu upływu; bez magnesu świeci szeroko całe okno dopiero od
        \SI{800}{V}. Wskazuje to na ogniskujące działanie pola
        magnetycznego.
  \item W~sesji z~magnesem prąd nasyca się przy ok.
        \SI{200}{\micro A} powyżej \SI{1200}{V}, co potwierdza również
        przestający rosnąć obraz plamki.
  \item Jasność ekranu silnie koreluje z~prądem emisji
        ($r_\mathrm{P}=0{,}95$--$0{,}99$) i~może służyć jako jego względna miara.
  \item Bezpośrednie porównanie wartości prądów między sesjami jest
        ograniczone: źródło było odgazowane tylko przed sesją bez magnesu,
        sesje różnią się napięciem $\UT$ (\SI{2000}{V} vs
        \SI{1500}{V}). Z~magnesem prąd jest 10--20 razy większy, co może
        wynikać zarówno z~ogniskowania, jak i~ze stanu źródła. Prąd
        upływu, nasycenie i~poświata w~sesji z~magnesem wynikają
        prawdopodobnie z~braku odgazowania.
  \item Do ilościowej oceny średnicy plamki potrzebne są: stały kadr
        i~ekspozycja kamery oraz wzorzec skali w~polu widzenia.
\end{enumerate}


\section{Widma EDS próbek}
\label{sec:widma}

Po potwierdzeniu, że emiter wytwarza wiązkę elektronów, zarejestrowano
widma promieniowania X wzbudzanego tą wiązką w~próbkach. Najpierw
zmierzono widmo próbki aluminiowej przy $\UT = \SI{4}{\kilo\volt}$,
a~następnie widma czterech próbek o~różnym składzie przy
$\UT = \SI{10}{\kilo\volt}$.
% Źródło: widma-eds/widma_tresc.tex (widmo 4 kV z 12.06.2026)
% =====================================================================
%  Widma EDS — fragment rozdziału 5 (sekcja \ref{sec:widma}).
%  Wstawiany we wspolne/rozdzialy/05_badania.tex po \section{Widma EDS próbek}.
%  Dane: widma-eds/dane_mca/2026-06-12_Al_4kV/4SERIA.mca (wykres: skrypty_mca/widma_mca.py).
%  Rysunki: widma-eds/widmo_eds_4kV.pdf, widma-eds/stabilnosc_zliczen.pdf
% =====================================================================

\subsection{Widmo próbki aluminiowej przy napięciu 4~kV}
\label{sec:widma-warunki}

Widmo zarejestrowano 12.06.2026 przy napięciu przyspieszającym
$\UT=\SI{4000}{\volt}$ i~napięciu ekstrakcyjnym $\UB=\SI{1000}{\volt}$
spektrometrem opisanym w~\pref{sec:rejestracja-widm}. Ustawienia toru
akwizycji zestawiono w~\tabref{tab:widma-parametry}.

\begin{table}[htbp]\centering
  \caption{Parametry rejestracji widma EDS}\label{tab:widma-parametry}
  \begin{tabular}{@{}ll@{}}\toprule
    Parametr & Wartość \\\midrule
    Napięcie przyspieszające $\UT$ & \SI{4000}{\volt} \\
    Napięcie ekstrakcyjne $\UB$ & \SI{1000}{\volt} \\
    Czas kształtowania (\angielski{peaking time}) & \SI{0.5}{\micro\second} \\
    Liczba kanałów, szerokość kanału & 8192, \SI{10}{\electronvolt} \\
    Czas rzeczywisty / czas życia & \SI{210.7}{\second} / \SI{100.1}{\second} \\
    Częstość zliczeń wejściowa / wyjściowa & ok.~475~kcps / ok.~225~kcps \\
    Czas martwy & \SIrange{52}{54}{\percent} \\\bottomrule
  \end{tabular}
\end{table}

Widmo zbierano w~trybie ciągłym z~kumulacją, a~stan pomiaru zapisano
w~ośmiu chwilach pośrednich. Badaną próbką była płytka aluminiowa, a~pomiar
wykonano z~magnesem ogniskującym umieszczonym na emiterze. Ciśnienie
w~komorze nie przekraczało \SI{1e-5}{\milli\bar}.
\TODO{prąd emisji $\Ie$ podczas pomiaru, jeśli zapisano}

\label{sec:widma-linie}%
Widmo po \SI{100}{\second} czasu życia przedstawiono na \rysref{fig:widmo-4kV}.
Skalę energii pliku programu ProSpect skorygowano liniowo na podstawie
położeń linii Mg~\Ka, Al~\Ka{} i~Si~\Ka{} (\pref{sec:widma-probki}).

\rysunek[0.95\linewidth]{widma-eds/widmo_eds_4kV.pdf}{Widmo EDS zarejestrowane
  przy $\UT=\SI{4}{\kilo\volt}$ i~$\UB=\SI{1}{\kilo\volt}$, czas życia \SI{100}{\second};
  w~powiększeniu zakres pików sumacyjnych \opracowanieWlasne}{fig:widmo-4kV}

W~widmie występują trzy linie charakterystyczne lekkich pierwiastków
(\tabref{tab:widma-linie}). Najsilniejsza jest linia Al~\Ka{} (ok.~$1{,}33\cdot10^{6}$
zliczeń w~kanale \SI{1.49}{\kilo\electronvolt}), a~linie Si~\Ka{} i~Mg~\Ka{} mają
wysokość odpowiednio ok.~\SI{63}{\percent} i~\SI{39}{\percent} piku Al.
Szerokość połówkowa pików wynosi ok.~\SI{87}{\electronvolt}.

\begin{table}[htbp]\centering
  \caption{Linie zidentyfikowane w~widmie z~\rysref{fig:widmo-4kV}}\label{tab:widma-linie}
  \begin{tabular}{@{}S[table-format=1.2]lS[table-format=1.3]l@{}}\toprule
    {$E$ piku [\si{\kilo\electronvolt}]} & Linia & {$E$ tabel. [\si{\kilo\electronvolt}]} & Pochodzenie \\\midrule
    1.26 & Mg~\Ka & 1.254 & domieszka w~próbce Al \\
    1.48 & Al~\Ka & 1.487 & próbka \\
    1.74 & Si~\Ka & 1.740 & struktura krzemowa\,/\,domieszka w~próbce \\
    2.98 & Al~\Ka{} + Al~\Ka & 2.974 & pik sumacyjny \\
    3.22 & Al~\Ka{} + Si~\Ka & 3.227 & pik sumacyjny \\
    3.46 & Si~\Ka{} + Si~\Ka & 3.480 & pik sumacyjny \\\bottomrule
  \end{tabular}
\end{table}

Ciągłe tło promieniowania hamowania zanika przy energii ok.~\SI{4}{\kilo\electronvolt},
co odpowiada granicy Duane'a--Hunta $E_\mathrm{max}=e\UT$ (\pref{sec:rtg}).
Elektrony docierają więc do próbki z~pełną energią ok.~\SI{4}{\kilo\electronvolt},
a~napięcie bramki nie zwiększa energii wiązki. Przy tej energii współczynnik
nadnapięcia wynosi ok.~3,1 dla Mg~K, 2,6 dla Al~K i~2,2 dla Si~K, co zapewnia
skuteczne wzbudzenie tych linii; linii K wapnia i~pierwiastków cięższych
nie można wzbudzić (\pref{subsec:jonizacja}).

\label{sec:widma-stabilnosc}%
W~widmie widoczne są także słabe piki przy \SIlist{2.98;3.22;3.46}{\kilo\electronvolt}, odpowiadające sumom
energii par fotonów Al i~Si zarejestrowanych w~jednym impulsie. Ich obecność
wynika z~bardzo dużej częstości zliczeń, przy której czas martwy przekracza
\SI{50}{\percent}. W~kolejnych pomiarach należy zmniejszyć prąd wiązki lub
zwiększyć odległość detektora od próbki, tak aby czas martwy nie przekraczał
ok.~\SI{30}{\percent}.

Przebieg częstości zliczeń w~czasie pomiaru przedstawiono na \rysref{fig:widma-stabilnosc}.
Częstości wejściowa i~wyjściowa pozostają stałe w~granicach \SI{\pm2}{\percent},
natomiast średnia częstość zliczeń w~kanale Al~\Ka{} maleje w~ciągu \SI{200}{\second}
z~ok.~$15{,}0\cdot10^{3}$ do $13{,}3\cdot10^{3}$~\si{\per\second}, tj.~o~ok.~\SI{11}{\percent}.
Spadek ten może wynikać ze zmian prądu emisji (podrozdz.~\ref{sec:emisja-magnes}), jednak
przy stałej częstości wejściowej bliskiej nasycenia toru akwizycji nie można
tego rozstrzygnąć bez pomiaru $\Ie$.

\rysunek[0.9\linewidth]{widma-eds/stabilnosc_zliczen.pdf}{Częstość zliczeń wejściowa (ICR),
  wyjściowa (OCR) i~średnia częstość w~kanale Al~\Ka{} w~czasie rejestracji widma
  \opracowanieWlasne}{fig:widma-stabilnosc}

Uzyskane widmo potwierdza, że wiązka elektronów z~emitera CNT wzbudza
w~próbce charakterystyczne promieniowanie X, a~linia Al~\Ka{} wskazana
w~celu pracy jest rejestrowana z~dużą liczbą zliczeń już przy
$\UT=\SI{4}{\kilo\volt}$.

% Widma próbek A–D z 13.11.2026 i linia tlenu: widma-eds/widma_materialy_tresc.tex
% =====================================================================
%  Widma EDS próbek o różnym składzie (pomiary 13.11.2026) — dalsza część
%  sekcji \ref{sec:widma}. Wstawiany we wspolne/rozdzialy/05_badania.tex
%  po % =====================================================================
%  Widma EDS — fragment rozdziału 5 (sekcja \ref{sec:widma}).
%  Wstawiany we wspolne/rozdzialy/05_badania.tex po \section{Widma EDS próbek}.
%  Dane: widma-eds/dane_mca/2026-06-12_Al_4kV/4SERIA.mca (wykres: skrypty_mca/widma_mca.py).
%  Rysunki: widma-eds/widmo_eds_4kV.pdf, widma-eds/stabilnosc_zliczen.pdf
% =====================================================================

\subsection{Widmo próbki aluminiowej przy napięciu 4~kV}
\label{sec:widma-warunki}

Widmo zarejestrowano 12.06.2026 przy napięciu przyspieszającym
$\UT=\SI{4000}{\volt}$ i~napięciu ekstrakcyjnym $\UB=\SI{1000}{\volt}$
spektrometrem opisanym w~\pref{sec:rejestracja-widm}. Ustawienia toru
akwizycji zestawiono w~\tabref{tab:widma-parametry}.

\begin{table}[htbp]\centering
  \caption{Parametry rejestracji widma EDS}\label{tab:widma-parametry}
  \begin{tabular}{@{}ll@{}}\toprule
    Parametr & Wartość \\\midrule
    Napięcie przyspieszające $\UT$ & \SI{4000}{\volt} \\
    Napięcie ekstrakcyjne $\UB$ & \SI{1000}{\volt} \\
    Czas kształtowania (\angielski{peaking time}) & \SI{0.5}{\micro\second} \\
    Liczba kanałów, szerokość kanału & 8192, \SI{10}{\electronvolt} \\
    Czas rzeczywisty / czas życia & \SI{210.7}{\second} / \SI{100.1}{\second} \\
    Częstość zliczeń wejściowa / wyjściowa & ok.~475~kcps / ok.~225~kcps \\
    Czas martwy & \SIrange{52}{54}{\percent} \\\bottomrule
  \end{tabular}
\end{table}

Widmo zbierano w~trybie ciągłym z~kumulacją, a~stan pomiaru zapisano
w~ośmiu chwilach pośrednich. Badaną próbką była płytka aluminiowa, a~pomiar
wykonano z~magnesem ogniskującym umieszczonym na emiterze. Ciśnienie
w~komorze nie przekraczało \SI{1e-5}{\milli\bar}.
\TODO{prąd emisji $\Ie$ podczas pomiaru, jeśli zapisano}

\label{sec:widma-linie}%
Widmo po \SI{100}{\second} czasu życia przedstawiono na \rysref{fig:widmo-4kV}.
Skalę energii pliku programu ProSpect skorygowano liniowo na podstawie
położeń linii Mg~\Ka, Al~\Ka{} i~Si~\Ka{} (\pref{sec:widma-probki}).

\rysunek[0.95\linewidth]{widma-eds/widmo_eds_4kV.pdf}{Widmo EDS zarejestrowane
  przy $\UT=\SI{4}{\kilo\volt}$ i~$\UB=\SI{1}{\kilo\volt}$, czas życia \SI{100}{\second};
  w~powiększeniu zakres pików sumacyjnych \opracowanieWlasne}{fig:widmo-4kV}

W~widmie występują trzy linie charakterystyczne lekkich pierwiastków
(\tabref{tab:widma-linie}). Najsilniejsza jest linia Al~\Ka{} (ok.~$1{,}33\cdot10^{6}$
zliczeń w~kanale \SI{1.49}{\kilo\electronvolt}), a~linie Si~\Ka{} i~Mg~\Ka{} mają
wysokość odpowiednio ok.~\SI{63}{\percent} i~\SI{39}{\percent} piku Al.
Szerokość połówkowa pików wynosi ok.~\SI{87}{\electronvolt}.

\begin{table}[htbp]\centering
  \caption{Linie zidentyfikowane w~widmie z~\rysref{fig:widmo-4kV}}\label{tab:widma-linie}
  \begin{tabular}{@{}S[table-format=1.2]lS[table-format=1.3]l@{}}\toprule
    {$E$ piku [\si{\kilo\electronvolt}]} & Linia & {$E$ tabel. [\si{\kilo\electronvolt}]} & Pochodzenie \\\midrule
    1.26 & Mg~\Ka & 1.254 & domieszka w~próbce Al \\
    1.48 & Al~\Ka & 1.487 & próbka \\
    1.74 & Si~\Ka & 1.740 & struktura krzemowa\,/\,domieszka w~próbce \\
    2.98 & Al~\Ka{} + Al~\Ka & 2.974 & pik sumacyjny \\
    3.22 & Al~\Ka{} + Si~\Ka & 3.227 & pik sumacyjny \\
    3.46 & Si~\Ka{} + Si~\Ka & 3.480 & pik sumacyjny \\\bottomrule
  \end{tabular}
\end{table}

Ciągłe tło promieniowania hamowania zanika przy energii ok.~\SI{4}{\kilo\electronvolt},
co odpowiada granicy Duane'a--Hunta $E_\mathrm{max}=e\UT$ (\pref{sec:rtg}).
Elektrony docierają więc do próbki z~pełną energią ok.~\SI{4}{\kilo\electronvolt},
a~napięcie bramki nie zwiększa energii wiązki. Przy tej energii współczynnik
nadnapięcia wynosi ok.~3,1 dla Mg~K, 2,6 dla Al~K i~2,2 dla Si~K, co zapewnia
skuteczne wzbudzenie tych linii; linii K wapnia i~pierwiastków cięższych
nie można wzbudzić (\pref{subsec:jonizacja}).

\label{sec:widma-stabilnosc}%
W~widmie widoczne są także słabe piki przy \SIlist{2.98;3.22;3.46}{\kilo\electronvolt}, odpowiadające sumom
energii par fotonów Al i~Si zarejestrowanych w~jednym impulsie. Ich obecność
wynika z~bardzo dużej częstości zliczeń, przy której czas martwy przekracza
\SI{50}{\percent}. W~kolejnych pomiarach należy zmniejszyć prąd wiązki lub
zwiększyć odległość detektora od próbki, tak aby czas martwy nie przekraczał
ok.~\SI{30}{\percent}.

Przebieg częstości zliczeń w~czasie pomiaru przedstawiono na \rysref{fig:widma-stabilnosc}.
Częstości wejściowa i~wyjściowa pozostają stałe w~granicach \SI{\pm2}{\percent},
natomiast średnia częstość zliczeń w~kanale Al~\Ka{} maleje w~ciągu \SI{200}{\second}
z~ok.~$15{,}0\cdot10^{3}$ do $13{,}3\cdot10^{3}$~\si{\per\second}, tj.~o~ok.~\SI{11}{\percent}.
Spadek ten może wynikać ze zmian prądu emisji (podrozdz.~\ref{sec:emisja-magnes}), jednak
przy stałej częstości wejściowej bliskiej nasycenia toru akwizycji nie można
tego rozstrzygnąć bez pomiaru $\Ie$.

\rysunek[0.9\linewidth]{widma-eds/stabilnosc_zliczen.pdf}{Częstość zliczeń wejściowa (ICR),
  wyjściowa (OCR) i~średnia częstość w~kanale Al~\Ka{} w~czasie rejestracji widma
  \opracowanieWlasne}{fig:widma-stabilnosc}

Uzyskane widmo potwierdza, że wiązka elektronów z~emitera CNT wzbudza
w~próbce charakterystyczne promieniowanie X, a~linia Al~\Ka{} wskazana
w~celu pracy jest rejestrowana z~dużą liczbą zliczeń już przy
$\UT=\SI{4}{\kilo\volt}$.
.
%  Dane: widma-eds/dane_mca/2026-11-13/ (.mca), skrypt: skrypty_mca/widma_mca.py,
%  położenia pików: piki_13-11.csv, rysunek: widma_materialy.pdf.
% =====================================================================

\subsection{Widma próbek o~różnym składzie}
\label{sec:widma-probki}

Kolejne widma zarejestrowano 13.11.2026 dla czterech próbek: A~(miedź o~nieznanym
składzie i~zanieczyszczonej powierzchni), B~(ołów), C~(złoto, chrom i~miedź) i~D~(aluminium).
Dla próbek A--C wykonano po 6--10 serii o~czasie życia
\SIrange{40}{250}{\second}, które zsumowano. Dla próbki~D wykorzystano serię
o~najmniejszym czasie martwym (\SI{20}{\percent}); pozostałe serie miały czas
martwy do \SI{90}{\percent}, przy którym widmo zniekształcają piki sumacyjne. Napięcie przyspieszające wynosiło
$\UT=\SI{10}{\kilo\volt}$, co potwierdza koniec tła promieniowania hamowania
przy ok.~\SI{10}{\kilo\electronvolt} (granica Duane'a--Hunta), więc
$\Ezero\approx\SI{10}{\kilo\electronvolt}$.

Skala energii zapisana przez program ProSpect (\num{8192} kanały,
\SI{80}{\kilo\electronvolt}) wymagała korekcji: linia Cu~\Ka{} wypadała
w~niej przy \SI{7.89}{\kilo\electronvolt} zamiast \SI{8.04}{\kilo\electronvolt}.
Zależność liniową $E = a\,E_\mathrm{p} + b$ wyznaczono z~położeń pięciu linii
o~jednoznacznej identyfikacji (Al~\Ka, Pb~M$\alpha$, Cr~\Ka, Fe~\Ka, Cu~\Ka).
Otrzymano $a = \num{1.0216}$ i~$b = \SI{-18}{\electronvolt}$, a~odchyłki
punktów kalibracyjnych nie przekraczają \SI{3}{\electronvolt}.
Pozostałe linie po korekcji leżą w~odległości do \SI{12}{\electronvolt}
od wartości tabelarycznych~\cite{goldstein2018sem}.

\rysunek[0.95\linewidth]{widma-eds/widma_materialy.pdf}{Widma EDS próbek A, B i~C
  (pomiary z~13.11.2026, skala logarytmiczna, zliczenia na sekundę czasu życia);
  szare pole -- zakres energii poza czułością toru pomiarowego, linia przerywana --
  położenie linii O~\Ka{} \opracowanieWlasne}{fig:widma-probki}

\begin{table}[htbp]\centering\small
  \caption{Linie zidentyfikowane w~widmach z~13.11.2026}\label{tab:widma-probki}
  \begin{tabular}{@{}lp{0.47\linewidth}p{0.3\linewidth}@{}}\toprule
    Próbka & Linie (energia po korekcji, \si{\kilo\electronvolt}) & Interpretacja \\\midrule
    A & Cu~L$\alpha$ 0,92; Cu~\Ka{} 8,04; Cu~\Kb{} 8,90; Te~L$\alpha$ 3,76,
        L$\beta_1$ 4,04, L$\beta_2$ 4,30, L$\gamma_1$ 4,57; P~\Ka{} 2,01;
        S~\Ka{} 2,30; Cl~\Ka{} 2,62; K~\Ka{} 3,32; Si~\Ka{} 1,73; Al~\Ka{} 1,48
      & miedź; Te i~P (prawdopodobnie dodatki stopowe);
        S, Cl, K, Si, Al z~zanieczyszczeń powierzchni \\
    B & Pb~M$\zeta$ 1,83; Pb~M$\alpha$ 2,35; Pb~M$\gamma$ 2,65 & ołów \\
    C & Au~M$\zeta$ 1,66; Au~M$\alpha\beta$ 2,13; Au~M$\gamma$ 2,41;
        Cr~\Ka{} 5,42; Cr~\Kb{} 5,95; Cu~L$\alpha$ 0,93; Cu~\Ka{} 8,04
      & złoto, chrom, miedź \\
    D & Al~\Ka{} 1,48; Mn~\Ka{} 5,91; Fe~\Ka{} 6,40; pik sumacyjny 2,98
      & stop aluminium z~Mn i~Fe \\\bottomrule
  \end{tabular}
\end{table}

Widma przedstawiono na \rysref{fig:widma-probki}, a~zidentyfikowane linie
zestawiono w~\tabref{tab:widma-probki}. W~próbce~A dominuje linia Cu~L$\alpha$.
Linie Cu~K są słabe, ponieważ przy $\Ezero\approx\SI{10}{\kilo\electronvolt}$
współczynnik nadnapięcia dla krawędzi Cu~K (\SI{8.98}{\kilo\electronvolt})
wynosi tylko ok.~1,1. Grupę pików przy \SIrange{3.76}{4.57}{\kilo\electronvolt}
przypisano serii L telluru: cztery linie zgadzają się z~wartościami tabelarycznymi
w~granicach \SI{9}{\electronvolt}, a~stosunek wysokości L$\beta_1$/L$\alpha$
(ok.~0,4) wyklucza alternatywną interpretację jako Ca~\Ka{} i~Ca~\Kb{}
(\SI{3.69}{} i~\SI{4.01}{\kilo\electronvolt}). Skład próbki nie jest znany. Tellur i~fosfor są typowymi
dodatkami miedzi tellurowej (CuTeP), a~S, Cl, K, Si i~Al pochodzą najpewniej
z~zanieczyszczeń powierzchni.
W~próbce~B widoczne są wyłącznie linie M ołowiu; linie Pb~L nie występują,
ponieważ krawędź Pb~L$_3$ (\SI{13.0}{\kilo\electronvolt}) leży powyżej $\Ezero$.
Widmo próbki~C zawiera linie M złota oraz linie Cr~K i~Cu, co odpowiada
warstwie złota na podłożu miedzianym z~chromem jako warstwą pośrednią.
Częstość zliczeń w~zakresie linii Au~M$\alpha$ w~dziesięciu kolejnych seriach
wynosiła \SIrange{1027}{1096}{\per\second}, tj.~zmieniała się w~granicach
ok.~\SI{\pm3}{\percent}. Widmo próbki~D jest podobne do \rysref{fig:widmo-4kV},
lecz nie zawiera linii Mg i~Si, a~zawiera słabe linie Mn i~Fe, więc jest to
inny stop aluminium.

\label{sec:widma-tlen}%
Osobnego omówienia wymaga zakres niskich energii, w~którym leży linia
tlenu. We wszystkich widmach liczba zliczeń spada poniżej ok.~\SI{0.65}{\kilo\electronvolt}
o~kilka rzędów wielkości, praktycznie do zera
(szare pole na \rysref{fig:widma-probki}). Linia O~\Ka{} (\SI{0.525}{\kilo\electronvolt})
leży poniżej tej granicy, dlatego tlen nie jest rejestrowany, mimo że powierzchnia
każdej z~próbek metalicznych jest pokryta tlenkiem rodzimym (Cu$_2$O/CuO, PbO, Al$_2$O$_3$)
i~zaadsorbowaną wodą. Przy ciśnieniu \SI{1e-5}{\milli\bar} warstwa adsorbatu
odnawia się w~czasie rzędu sekundy, więc pomiar w~komorze nie usuwa tlenu
z~powierzchni. Dolną granicę wyznacza najpewniej absorpcja fotonów w~oknie wejściowym
i~warstwie nieczułej detektora oraz próg dyskryminacji impulsów toru
akwizycji~\cite{goldstein2018sem}; typu okna użytego detektora nie ustalono.

Najsilniejszy pik w~zakresie niskich energii, przy \SI{0.92}{\kilo\electronvolt},
nie pochodzi od tlenu, lecz jest linią Cu~L$\alpha$ (\SI{0.930}{\kilo\electronvolt}).
Potwierdzają to linie Cu~\Ka{} i~Cu~\Kb{} w~tym samym widmie, obecność tego piku
w~próbce~C z~podłożem miedzianym i~jego brak w~próbkach B i~D. Odległość od
linii O~\Ka{} wynosi ok.~\SI{400}{\electronvolt}, tj.~pięciokrotność szerokości
połówkowej piku (ok.~\SI{80}{\electronvolt}), więc nakładanie obu linii jest
wykluczone. Pik leży tuż nad granicą czułości, przez co jego wysokość względem
linii o~większych energiach jest zaniżona.

Do oznaczania tlenu, najczęstszego pierwiastka skał i~regolitu, analizator
docelowy wymaga detektora z~cienkim oknem polimerowym lub bez okna.


Wyniki obu części badań zestawiono i~odniesiono do celu pracy
w~\rozdzref{ch:podsumowanie}, gdzie wskazano także ograniczenia
prototypu i~kierunki dalszych prac.
